The mass-conservation finding (Section~\ref{sec:methods}, Figure~\ref{fig:e02}) is worth dwelling
on because it is a natural but wrong instinct: a modeler who checks $S+I+R \approx N$ as a sanity
test on their SIR solver will see it satisfied to roundoff \emph{even for a badly wrong solve},
because the property follows from the right-hand side summing to zero identically, not from
accurate integration. The phase invariant $Q(t)$ derived from the same exact reduction that gives
our gold standard is the diagnostic that actually carries information, and it is available with no
extra computational cost.

The crossover result (Figure~\ref{fig:e13}) reframes the base paper's implicit assumption. Capaldi
et al.\ do not need to specify a solver precisely because, in the noise regime they study, sampling
uncertainty already dominates; our contribution is to make that dominance relationship explicit and
locate where it stops holding. This is a genuinely two-sided result: it is not simply ``better
solvers are always better'' (uninteresting) or ``solver choice never matters'' (also uninteresting,
and false at low noise) but a quantitative boundary between two regimes, each of which is the
correct practical conclusion in its own part of the parameter space.
